% ===========================================================================
% multitask_bo.tex -- the running technical report.
%
% This document keeps every experiment in the order it was run, including the
% ones that failed and the two conclusions the null control reversed. It is
% deliberately NOT reorganised: the order is part of the evidence, because
% several claims here are only interpretable next to the check that came after
% them.
%
% The same work rewritten for a reader who has not followed the project --
% abstract, introduction, related work, overview, method, setup, results,
% discussion, conclusion -- is report/paper.tex. That document also carries the
% later KIBA material (history-seeded initial designs, the coupling ladder, the
% bad-start comparison) which is not in this one.
% ===========================================================================

\documentclass[11pt]{article}

\usepackage[margin=1in]{geometry}
\usepackage{amsmath,amssymb}
\usepackage{booktabs}
\usepackage{graphicx}
\usepackage{xcolor}
\usepackage[hidelinks]{hyperref}
\usepackage{enumitem}
\usepackage{caption}

\captionsetup{font=small}
\newcommand{\stl}{\textsc{stl}}
\newcommand{\mt}{\textsc{mt}}
\newcommand{\pooled}{\textsc{pooled}}
\newcommand{\simarm}{\textsc{sim}}
\newcommand{\rnd}{\textsc{rand}}

\title{When Does Sharing Data Between Related Targets Help\\
       Machine-Guided Compound Selection?}
\author{}
\date{}

\begin{document}
\maketitle

\begin{abstract}
\noindent
Testing compounds is expensive, so programmes test a few at a time and use a
model to choose the next few. A new protein target starts with no data, while
related targets often have years of it, so sharing data between them is an
obvious thing to try. We ran that comparison carefully on two measured protein
panels and found that the usual way of running it can produce a gain that is not
there.

\paragraph{The check.} Give the sharing model an extra data column that is
\textbf{pure noise}. It carries no information, so any systematic gain cannot be information
transfer. It produced
up to $\mathbf{+14.7}$ points of extra recovery of the best compounds. The predicted values do not explain the gain: the model correctly judges the
noise column to be weakly related and its predictions get slightly worse. The cause is not the extra data as such. The model's signal scale is estimated
across all columns at once, and a weak between-target correlation can be fitted
from a finite sample, so adding an uninformative column changes how much
uncertainty the model reports even when its predictions are unchanged to
fourteen decimal places. Expected improvement reads that uncertainty, and
therefore picks different compounds. This check reversed two of our own
conclusions, one of which had already passed five other checks.

\paragraph{It depends on the picking rule.} Of the three rules we tried, only
expected improvement lets a noise column \emph{gain} recovery. Under the other two the null effect was neutral or negative rather than
spuriously positive. Both directions are distortions of a comparison that should
show no systematic informational benefit; only a positive one manufactures
apparent transfer. The
real effect behaves the same way: on one panel, sharing data is worth $+7.2$
points under one rule, $+16.2$ under another, and nothing at all under a third.
No two rules agree on both panels about whether sharing helps.

\paragraph{What is actually worth doing.} Interpreted against the matched null-auxiliary control,
sharing data between targets is worth $+7$ to $+14$ points on these two
panels, when the old data was already paid for and was measured on the same
compounds. The gains were largest and most consistent around $22$--$33\%$
library coverage. In one campaign-matching calculation, roughly five historical
measurements on related targets were needed to replace one measurement on the
target of interest, so the result is most attractive when that data already
exists. By
contrast, simply changing \emph{which number the model predicts} --- a
selectivity score instead of raw potency --- is worth $+39.7$ points at the same
cost on the same data, with no sharing at all. The choice of what to optimise produced one of the largest effects we measured.

We also list eight things we expected to be true that were not, and two data
problems in widely used public benchmarks. Everything here is measured by
replaying published data; no compound was made or tested for this work.
\end{abstract}

\section{Motivation}

A campaign against a new target begins with no data on it, a library of
available compounds, and a budget of a few hundred assays. BO fits that: fit a
surrogate to what is measured, pick the next batch by an acquisition function,
repeat. Targets rarely arrive alone --- a kinase programme usually holds years
of assay data on related kinases --- so modelling them jointly is an obvious
thing to try, and multi-task Gaussian processes are the obvious way to try it.

Evaluating that idea requires comparing two surrogates under a fixed
acquisition function, and this is where the problem lies. Expected improvement
reads the model's prediction \emph{and} the model's uncertainty, and that
uncertainty is not invariant to the auxiliary dataset. The signal scale is
estimated across all columns at once, so adding columns moves it in either
direction, and a weak between-target correlation can be fitted from a finite
sample. A variance-sensitive selection rule then behaves differently for
reasons unrelated to what the model has learned. On our
panels that effect alone is worth up to $14.7$ recall points --- larger than the
transfer effect being measured.

The control that detects this is cheap: give the multi-task arm an auxiliary
that is pure noise and test whether it ties the single-task arm. On our
panels it does not always tie, which is the result. We did not find it
in the multi-task molecular BO literature, we did not run it ourselves for most
of this project, and running it reversed two of our conclusions. We therefore
report it first (\S\ref{sec:null}) and interpret every later effect size
against the matched null-auxiliary control.

With the control in place, the benchmark still has a positive result, but not
the one we set out to find. Objective specification produced one of the
largest effects we measured, and in the selectivity experiment it mattered more
than joint versus single-target modelling. We organise the remainder around the decisions a campaign
makes, each with what it is worth and what it costs.

\section{Related work}
\label{sec:related}

\paragraph{Retrospective evaluation is the norm.} Chaudhry et al.~\cite{chaudhry2026practical} benchmark
foundation-model-driven BO across 61 regression tasks spanning proteins, DNA,
RNA and small molecules, entirely on established retrospective benchmarks. Dang
et al.\ (CheapVS) screen a 100{,}000-molecule library with a docking oracle and
validate by recovering known drugs --- 16/37 for EGFR and 37/58 for DRD2 ---
rather than by new measurement. We adopt the same frame and state its limits
explicitly (\S\ref{sec:limitations}).

\paragraph{Prospective work sets the calibration.} Where campaigns do reach the
bench, the reported quantity is hit-rate enrichment rather than top-$k$ recall.
ChemScreener~\cite{chemscreener2026} raised a WDR5 hit rate from $0.49\%$ (primary HTS) to $5.91\%$ over
five iterative HTRF rounds; a structure-based deep-learning pipeline with an
automated lab~\cite{irak1_2024} recovered $23.8\%$ of IRAK1 hits within the top $1\%$ of ranked
compounds; BATCHIE~\cite{batchie2023} performs Bayesian active learning for combination screens on
cell lines. Our single-task arm recovers $27$--$38\%$ of the top-10 on panels of
a few hundred compounds, which is not the same metric and not a comparable
number: our libraries are enriched in actives by construction.

\paragraph{The surrogate is moving away from GPs.} Chaudhry et al.~\cite{chaudhry2026practical} find GP
surrogates consistently outperformed by neural alternatives with Thompson
sampling, particularly MLP ensembles and variational Bayesian last layers, and
report that smaller batch sizes reach peak performance faster and recover more
elite candidates at equal total budget. We use an exact GP with EI throughout,
which bounds the scope of our surrogate comparisons; \S\ref{sec:null} reports
the noise control under Thompson sampling as well as EI, because the
confound originates in the surrogate's uncertainty calibration and manifests
through the acquisition, so its size depends on the interaction between the
two.

\paragraph{Objective specification over model choice.} CheapVS~\cite{dang2025cheapvs} obtains its
efficiency by eliciting chemist preferences over trade-offs among affinity,
molecular weight, lipophilicity and ADMET properties --- roughly 2{,}200
pairwise comparisons and about ten hours of expert time per experiment --- not
by a better surrogate. Our \S\ref{sec:selectivity} reaches the same conclusion
with a fixed contrast that requires no human input, and we regard the agreement
between the two as the stronger evidence.

\paragraph{Leakage in affinity predictors.} Reported accuracies of co-folding
affinity models on public benchmarks are contested~\cite{boltz2_prereview,leakage2026}: their training corpora
overlap the public binding databases these benchmarks are drawn from, and
splitting by protein-sequence identity is insufficient because homologous
targets retain correlated binding profiles. We therefore do not evaluate such a
model here (\S\ref{sec:next}).

\section{Setup}

Let $\mathcal{M}$ be a fixed library of $N$ existing compounds and
$\mathcal{T} = \{t_0, \ldots, t_{T-1}\}$ a set of protein targets, with $t_0$
the \emph{primary} target we wish to optimise. The oracle $y(m, t)$ returns a
measured potency (pIC$_{50}$, higher is better).

\paragraph{Budget is counted in cells.} A \emph{cell} is one
$(\text{compound}, \text{target})$ measurement. This accounting is not
incidental. A single-task campaign spends all $n$ cells on $t_0$; a multi-task
campaign that measures every target on each compound affords only $n/T$
compounds. Comparing the two at equal \emph{compounds} rather than equal cells
hands the multi-task arm more data and makes any result it produces
unfalsifiable. Several apparently positive multi-task results in our own early
runs disappeared once this was enforced.

\paragraph{Two hypotheses, not one.}
\begin{description}[leftmargin=2em,style=nextline]
  \item[budget-splitting] Given a fresh budget of $n$ cells, spreading them across
    $t_0$ and its relatives beats spending all $n$ on $t_0$.
  \item[borrowed-history] Given $n$ cells on $t_0$ \emph{plus} $H$ pre-existing
    observations per relative that cost this campaign nothing, joint modelling
    beats single-task modelling on $t_0$.
\end{description}
These differ in who paid for the auxiliary data. They have opposite answers.

\subsection{Data}

\begin{table}[ht]
\centering\small
\begin{tabular}{lrrl}
\toprule
Dataset & Compounds & Targets & Label \\
\midrule
DOCKSTRING            & 260{,}148 & 58 & AutoDock Vina score (simulated) \\
ChEMBL, separate libraries & 3{,}970  & 4 & pIC$_{50}$ (measured) \\
BindingDB JAK block   & 1{,}563 $\to$ 904 & 4 & pIC$_{50}$ (measured, co-measured) \\
BindingDB CYP block   & 1{,}472 $\to$ 335 & 4 & pIC$_{50}$ (measured, co-measured) \\
\bottomrule
\end{tabular}
\caption{Datasets. Arrows denote the compound count after the censoring filter
of \S\ref{sec:data-quality}. The ChEMBL libraries are per-target and share
$0.1$--$1.7\%$ of their compounds pairwise; the BindingDB blocks are fully
observed selectivity panels.}
\label{tab:data}
\end{table}

DOCKSTRING is dense by construction and lets us ask questions that measured data
cannot, at the cost of labels that are predicted rather than measured. The
ChEMBL libraries represent the common situation of separate, independently
assembled per-target datasets. The BindingDB blocks are \emph{selectivity
panels}: the same compounds measured against every member of a protein family,
which removes chemical-overlap confounds by construction.

Empirical cross-target correlations, measured on shared compounds, are
$+0.24$ to $+0.71$ (median $+0.57$) for the JAK block and $+0.43$ to $+0.68$
(median $+0.54$) for the CYP block --- strictly inside the range where transfer
is supposed to pay.

\paragraph{Two problems in the source data.}
\label{sec:data-quality}

These affect anyone using the same sources and are reported so they are not
rediscovered.

\paragraph{The published DOCKSTRING objectives are dominated by a free
descriptor.} Both objectives of Kubincov\'a \& Mobley carry a
$-10\,(1 - \mathrm{QED})$ term. QED is computed from the SMILES: it costs
nothing and is known for the whole library before any assay. That term has
$2.0$--$2.7\times$ the spread of the entire docking component, and
$\rho(\text{objective}, \text{QED term}) = +0.93$, while
$\rho(\text{objective}, -\text{JAK2}) = +0.007$. Ranking by QED alone, with zero
assays, recovers $48.1\%$ of the top-50; the best arm given 80 assays reaches
$50.6\%$. Any top-$k$ number on these objectives must be reported against that
floor. We therefore define potency variants with drug-likeness as a library
admission filter ($\mathrm{QED} \ge 0.5$) rather than a score term, which
reduces the free share from $95\%$ to $33\%$.

\paragraph{The BindingDB CYP block is $77\%$ censored.} $21$--$27\%$ of each CYP
column is exactly pIC$_{50} = 5.00$, with further spikes at $4.52$, $4.30$,
$5.30$, $4.00$ --- the fixed doses $10$, $30$, $50$, $5$, $100\,\mu$M of a
screening panel, stored as if they were potencies. $40.3\%$ of CYP compounds
carry the identical value on all four targets (JAK: $3.5\%$). This inflates
cross-target correlation, rails the fitted task correlation to 1, makes naive
pooling appear optimal, and turns ``transfer'' into the copying of a constant.
We filter compounds touching any value held by $>2\%$ of a column; $22.8\%$ of
the CYP block and $57.8\%$ of the JAK block survive. All measured results below
use the filtered blocks.

\subsection{Methods}

\paragraph{Surrogate.} Exact Gaussian process regression with a Tanimoto kernel
over 2048-bit ECFP4 fingerprints. Hyperparameters (signal variance, noise ratio)
are chosen by the model's own fit score over a fixed grid with an analytic
inner solve, so fits are deterministic given a seed.

\paragraph{The sharing model.} One kernel over compounds, scaled by a
between-target correlation $r$:
\[
  k\big((m, a), (m', b)\big) \;=\; k_{\mathrm{Tanimoto}}(m, m')\; B[a, b],
  \qquad
  B \;=\; (1 - r)\,I + r\,\mathbf{1}\mathbf{1}^{\!\top},
\]
with $r$ selected by the model's own fit score over a grid that includes both
endpoints. Targets are $z$-scored per task. $r = 0$ recovers independent
single-task GPs exactly; $r = 1$ is naive pooling. Each task carries its own
compound index, so blocked designs (all targets on the same compounds) and
spread designs (each target on its own compounds) use identical machinery.

\paragraph{Arms.}
\begin{description}[leftmargin=2em,style=nextline,itemsep=0.2em]
  \item[\rnd] Uniformly random selection. The floor.
  \item[\simarm] Similarity search: rank the library by maximum Tanimoto
    similarity to the most potent compounds known so far, from history and
    measurements alike. This is what a medicinal chemist does without a model,
    uses exactly the information \mt{} uses, and costs nothing.
  \item[\stl] Single-task Tanimoto GP + acquisition. Standard BO.
  \item[\pooled] The sharing model with $r$ forced to 1: the same observations
    with the task structure discarded. Separates ``modelling the task structure
    helped'' from ``more rows helped''.
  \item[\mt] The sharing model with $r$ fitted from the data. The method under test.
\end{description}

\paragraph{Controls on the comparison itself.} Two, both added after the
results below were first drafted and both of which changed them.

\begin{description}[leftmargin=2em,style=nextline,itemsep=0.25em]
  \item[Null-information auxiliary.] Each auxiliary column is replaced by a
    surrogate variable with controlled rank correlation $\rho$ to its own
    true values, mixing the $z$-scored column with independent noise. The
    primary target, the history \emph{indices} and the success set are
    untouched; only what the multi-task arm is \emph{told} about its relatives
    degrades. $\rho = 1$ reproduces the measured panel exactly (verified: the
    stored numbers are recovered to all printed digits), and $\rho = 0$ is the
    noise control. A fresh null column is drawn for each (target, repeat,
    auxiliary), giving 96 independent null realisations across the 32 campaigns
    in each reported configuration rather than reusing one. Their achieved
    correlations span $-0.12$ to $+0.12$ about zero.
  \item[Variance-free acquisition.] Expected improvement reads mean and
    variance. Ranking on the model's prediction alone (\texttt{greedy}) removes the
    variance channel, so an arm difference that survives it cannot be an
    artefact of variance calibration. \texttt{greedy} is a worse optimiser in
    absolute terms and we use it as a diagnostic, not as a proposed method.
\end{description}

\paragraph{Campaign protocol.} Identical initial design ($n_{\text{init}} = 10$),
budget, batch size and acquisition (expected improvement) across arms; only the
surrogate differs. Historical auxiliary observations are fixed before the
campaign and never grow --- they are history, not budget. Metrics: top-$k$
recall ($k = 10$) and AUC of the top-$k$ trace. Reported over $8$--$15$ repeats
$\times$ 4 held-out targets.

\subsection{Statistical analysis}
\label{sec:stats}

\paragraph{Unit of replication.} A cell is one (held-out target, repeat)
campaign. Each block has four held-out targets and we run eight repeats per
target, so a cell count of 32 is \emph{eight repeats nested within four
targets}, not 32 independent replicates. Repeats differ only in the random
draws --- the initial design and the free history --- so they are algorithmic,
not biological. Targets are the higher-level unit and there are four of them.

\paragraph{Pairing.} Arms within a (target, repeat) share the initial design and
the free history by construction, verified at runtime (\S\ref{sec:repro}).
Comparisons between arms are therefore paired at the campaign level and we use a
paired $t$-test on the per-campaign difference. Comparisons between
\emph{configurations} that do not share draws --- a random against a
scaffold-split seed, one library size against another --- are unpaired and use
Welch's $t$-test.

\paragraph{What these $p$-values do and do not establish.} Campaign-level
tests quantify algorithmic variability --- sensitivity to the randomised initial
design and history draw --- \emph{conditional on the four targets evaluated}.
They are not population-level inference over unseen protein targets. Eight seeds
on JAK1 do not make JAK1 eight targets. Generalisation across targets is
assessed descriptively, through the four held-out-target effects and their sign
consistency, because each block contains only four targets and no test can
manufacture a larger biological $n$ from that.

Accordingly we give the four target-level mean effects for every headline
result rather than only a ``4/4 up'' count:

\begin{center}\small
\begin{tabular}{llrrrr}
\toprule
Result & & \multicolumn{4}{c}{per-target effect (recall points)} \\
\midrule
Sharing, JAK $22\%$   & & $+11.2$ & $+8.8$  & $+3.8$  & $+2.5$ \\
Sharing, JAK $33\%$   & & $+15.0$ & $+7.5$  & $+5.0$  & $+1.2$ \\
Sharing, CYP $22\%$   & & $+8.8$  & $+17.5$ & $+15.0$ & $+16.2$ \\
Sharing, CYP $33\%$   & & $+11.3$ & $+12.5$ & $+5.0$  & $+18.7$ \\
Objective choice, JAK & & $+55.0$ & $+67.5$ & $+27.5$ & $+8.8$ \\
\bottomrule
\end{tabular}
\end{center}

The spreads are wide --- $+1.2$ to $+15.0$ within one JAK cell, $+8.8$ to
$+67.5$ for the objective result --- and that spread, not the pooled mean, is
what a practitioner facing a single target should read.

\paragraph{What we report.} The estimator is the mean per-campaign difference in
top-10 recall, in percentage points, with its standard error over the 32 cells.
We give $p$-values because reviewers expect them, but the quantity we ask the
reader to weigh is the effect size and its spread, and we report per-target
outcomes alongside every pooled figure. At four targets, a pooled mean can hide
opposite signs: in one budget cell one target gained 15 points while another
lost 15, pooling to a meaningless $+3.7\%$.

\paragraph{Multiplicity.} We make many comparisons --- two blocks, four targets
each, several budgets, three acquisition functions, several history
compositions, and a range of auxiliary qualities --- and we do \emph{not} apply
a family-wise correction. We therefore do not treat any single $p$ near $0.05$
as establishing a result. The claims we advance are the ones that hold with a
large effect across both blocks, or that reverse a conclusion, and we mark
borderline values as such rather than counting them as positive. Readers should
treat $p = 0.037$ in Table~\ref{tab:acqnull} as suggestive within a large
comparison set, not as a standalone finding; the claim it supports rests on the
$+14.7$ point effect on the other block.

\section{The noise control}
\label{sec:null}

We report this first because every later effect size is interpreted against the matched null-auxiliary control, and
because it reversed two conclusions that had passed every other check we
applied.

Every arm comparison above asks whether a surrogate given auxiliary data beats
one that is not. The control that makes this answerable is to give the
multi-task arm an auxiliary carrying no information and test whether it ties. We
found no multi-task molecular BO result that runs it.

\begin{table}[ht]
\centering\small
\begin{tabular}{llrrr}
\toprule
Block & Arm & $\rho = 0$ (no information) & $\rho = 1$ (measured) & Not transfer \\
\midrule
JAK & \mt{} $-$ \stl{}     & $+2.8$ \ ($p = 0.037$) & $+7.2$  & ${\sim}39\%$ \\
CYP & \mt{} $-$ \stl{}     & $+1.6$ \ (n.s.)        & $+11.9$ & ${\sim}13\%$ \\
CYP & \pooled{} $-$ \stl{} & $\mathbf{+14.7}$ \ ($p < 0.001$) & $+20.9$ & $\mathbf{{\sim}70\%}$ \\
\bottomrule
\end{tabular}
\caption{Top-10 recall points at $33\%$ of the library, top-10 withheld from
history, 32 campaigns per cell, expected improvement. A pure-noise auxiliary
carries no information, so any gain in the $\rho = 0$ column is not transfer.}
\label{tab:null}
\end{table}

\paragraph{The posterior mean does not explain it.} On the JAK block with 100 primary
and 150 auxiliary observations, the model's own fit score selects $r = 0.621$ on
measured auxiliaries and $r = 0.083$ on noise, and the noise posterior is worse
than single-task (rank correlation $0.715$ against $0.733$). The posterior degrades
while the campaign improves, so the gain does not come through the posterior
mean.

\paragraph{Where it becomes a different decision.} Mean predictive standard deviation is
$0.671$ (\stl{}), $0.644$ (\mt{}) and $0.561$ (\pooled{}). EI trades mean
against uncertainty, and adding auxiliary columns re-estimates the uncertainty
scale whether or not they are informative. On these blocks
that helps, and it is indistinguishable from transfer in the output. Ranking on
the model's prediction alone removes the channel:

\begin{table}[ht]
\centering\small
\begin{tabular}{llrrr}
\toprule
Block & Acquisition & $\rho$ & \mt{} $-$ \stl{} & \pooled{} $-$ \stl{} \\
\midrule
JAK & EI     & 0 & $+2.8$ \ ($p=0.037$) & $+1.9$ \ (n.s.) \\
JAK & EI     & 1 & $+7.2$ \ ($p<0.001$) & $+6.2$ \ ($p<0.001$) \\
JAK & greedy & 0 & $+3.1$ \ (\textbf{n.s.}) & $+7.5$ \ (n.s.) \\
JAK & greedy & 1 & $\mathbf{+16.2}$ \ ($p=0.008$) & $+4.7$ \ (n.s.) \\
\midrule
CYP & EI     & 0 & $+1.6$ \ (n.s.) & $\mathbf{+14.7}$ \ ($p<0.001$) \\
CYP & EI     & 1 & $+11.9$ \ ($p<0.001$) & $+20.9$ \ ($p<0.001$) \\
CYP & greedy & 0 & $-2.5$ \ (n.s.) & $\mathbf{-20.6}$ \ ($p=0.001$) \\
CYP & greedy & 1 & $+4.4$ \ (n.s.) & $\mathbf{-7.5}$ \ ($p=0.081$) \\
\bottomrule
\end{tabular}
\caption{The same comparisons with and without the variance channel.}
\label{tab:greedy}
\end{table}

\paragraph{Where the effect actually enters.} Our first explanation was that
extra kernel rows compress the model's uncertainty whether or not they are
informative. That is wrong in both direction and mechanism, and the correct
account is sharper.

With $r = 0$ the task covariance is the identity, the joint kernel is block
diagonal, and the auxiliary rows cannot touch the primary task at all --- not
its predictions and not its uncertainty. That is a mathematical statement, and
the first rung below confirms it numerically. But the signal variance is not a
free parameter: it is profiled out of the likelihood as
\[
  \hat{s} \;=\; \frac{1}{n}\,z^{\top}\!\left(K + \eta I\right)^{-1} z ,
\]
and for a joint fit both $z$ and $n$ run over \emph{every} task. $\hat{s}$ is
therefore a pooled quantity, and the predictive uncertainty
$\sigma = \sqrt{v\,\hat{s}}$ inherits it. Auxiliary rows change the primary
task's uncertainty through this one channel even when they are uncorrelated
with it.

We release one control at a time to confirm that this is the only channel.

\begin{table}[ht]
\centering\small
\begin{tabular}{lrrrrr}
\toprule
Rung & $\max_i|\Delta\mu_i|$ & $\max_i|\Delta\sigma_i|$ & mean $\sigma$
     & EI order vs \stl{} & $r$ \\
\midrule
\multicolumn{6}{l}{\emph{JAK}, single-task mean $\sigma = 0.6428$} \\
$r=0$, all pinned        & $1.2\!\times\!10^{-14}$ & $5.1\!\times\!10^{-14}$ & $0.6428$ & $1.000$ & $0.000$ \\
$r=0$, scale pooled      & $1.2\!\times\!10^{-14}$ & $\mathbf{12.5}$          & $1.2045$ & $\mathbf{0.181}$ & $0.000$ \\
$r=0$, scale and $\eta$  & $1.44$                  & $0.48$                   & $0.5582$ & $0.081$ & $0.000$ \\
$r$ fitted, noise aux.\  & $1.86$                  & $0.54$                   & $0.5549$ & $0.078$ & $0.159$ \\
$r$ fitted, real aux.\   & $1.92$                  & $0.43$                   & $0.6357$ & $0.068$ & $0.596$ \\
\midrule
\multicolumn{6}{l}{\emph{CYP}, single-task mean $\sigma = 0.6292$} \\
$r=0$, all pinned        & $5.3\!\times\!10^{-15}$ & $9.4\!\times\!10^{-15}$ & $0.6292$ & $1.000$ & $0.000$ \\
$r=0$, scale pooled      & $5.3\!\times\!10^{-15}$ & $\mathbf{5.31}$          & $0.7602$ & $\mathbf{0.842}$ & $0.000$ \\
$r=0$, scale and $\eta$  & $1.11$                  & $0.41$                   & $0.5215$ & $0.254$ & $0.000$ \\
$r$ fitted, noise aux.\  & $1.11$                  & $0.41$                   & $0.5194$ & $0.197$ & $0.151$ \\
$r$ fitted, real aux.\   & $1.88$                  & $0.49$                   & $0.5902$ & $0.094$ & $0.809$ \\
\bottomrule
\end{tabular}
\caption{Releasing one control at a time against the single-task fit. 24 draws
$\times$ 4 held-out targets, 100 primary observations, 50 per auxiliary. Rows
2--4 use auxiliary columns of pure noise. $\Delta$ is against the single-task
posterior on the same draw.}
\label{tab:mechanism}
\end{table}

\begin{enumerate}[leftmargin=1.6em,itemsep=0.25em]
  \item \textbf{The auxiliary rows are exactly inert when nothing is allowed to
    respond to them.} With $r = 0$ and the signal variance, noise ratio and
    normalisation all pinned to the single-task fit, predictions and uncertainty
    agree to machine precision ($10^{-14}$) and the acquisition ordering is
    unchanged. The block-diagonal argument holds, as it must.
  \item \textbf{Releasing the pooled signal variance alone changes the
    uncertainty, and that alone changes the decisions.} The predictions stay
    identical to $10^{-14}$, the uncertainty moves by up to $12.5$, and the
    expected-improvement ordering agreement falls to $0.181$ on JAK. Nothing
    about the model's estimate of the objective has changed; only the scale on
    which it reports its own uncertainty.
  \item \textbf{Finite-sample correlation is a second, smaller channel.} With
    $r$ free the fit selects $r = 0.159$ (JAK) and $0.151$ (CYP) on columns that
    are pure noise, and the predictions then degrade. Real auxiliaries move it
    the other way at $r = 0.596$ and $0.809$.
\end{enumerate}

The attribution confound is therefore specific and checkable: \emph{auxiliary
rows cannot affect the primary model when its parameterisation is held fixed;
the confound enters when the pooled scale estimate and the task correlation are
allowed to respond to the auxiliary dataset.} The acquisition function is where
that difference becomes a different decision, not where it originates.

\paragraph{The sign of the artefact depends on the acquisition.} Greedy removes
the variance channel entirely. Thompson sampling keeps it but uses it
differently: it draws $\mu + \sigma\varepsilon$, so a compressed $\sigma$ reduces
exploration noise rather than sharpening a mean/variance trade-off. Running the
noise control under all three separates them.

\begin{table}[ht]
\centering\small
\begin{tabular}{llrrr}
\toprule
Block & Acquisition & \stl{} & \mt{} $-$ \stl{} & \pooled{} $-$ \stl{} \\
\midrule
JAK & EI     & $76.2\%$ & $\mathbf{+2.8}$ \ ($p=0.037$) & $+1.9$ \ (n.s.) \\
JAK & greedy & $51.6\%$ & $+3.1$ \ (n.s.) & $+7.5$ \ (n.s.) \\
JAK & TS     & $78.4\%$ & $-3.4$ \ (n.s.) & $\mathbf{-8.1}$ \ ($p=0.001$) \\
\midrule
CYP & EI     & $58.1\%$ & $+1.6$ \ (n.s.) & $\mathbf{+14.7}$ \ ($p<0.001$) \\
CYP & greedy & $55.9\%$ & $-2.5$ \ (n.s.) & $\mathbf{-20.6}$ \ ($p=0.001$) \\
CYP & TS     & $48.8\%$ & $-1.2$ \ (n.s.) & $-0.9$ \ (n.s.) \\
\bottomrule
\end{tabular}
\caption{The noise control ($\rho = 0$) under three acquisition
functions. A noise auxiliary carries nothing, so every entry should be zero.
32 campaigns per cell.}
\label{tab:acqnull}
\end{table}

\textbf{Among the three acquisition functions we tested, expected improvement
was the only one showing a clear positive null-auxiliary deviation} ($+2.8$ on JAK, $+14.7$ for
\pooled{} on CYP). Under greedy and Thompson sampling the significant deviations
are negative --- a useless auxiliary degrades the posterior, and those
acquisitions register the degradation. Both directions are distortions of a
comparison that should read zero, but only EI manufactured apparent transfer in these
experiments, and only EI turned a null result into a positive one here.

\paragraph{The transfer effect is itself acquisition-dependent.} Measured with
the real panel ($\rho = 1$), \mt{} $-$ \stl{} is:

\begin{center}\small
\begin{tabular}{lrrr}
\toprule
Block & EI & greedy & TS \\
\midrule
JAK & $+7.2$ \ ($p<0.001$) & $+16.2$ \ ($p=0.008$) & $+0.6$ \ (n.s.) \\
CYP & $+11.9$ \ ($p<0.001$) & $+4.4$ \ (n.s.) & $+13.8$ \ ($p=0.003$) \\
\bottomrule
\end{tabular}
\end{center}

\textbf{On JAK the gain does not survive Thompson sampling} ($+0.6$, $p=0.86$),
having been $+7.2$ under EI and $+16.2$ under greedy. On CYP it survives TS
($+13.8$) but not greedy ($+4.4$, n.s.). No two acquisitions agree on both
blocks about whether joint modelling helps.

We take this as the strongest statement the data supports: \emph{multi-task
transfer is not a property of the surrogate that a benchmark can measure once.}
It is a property of the surrogate, the acquisition function and the panel
together, and a result reported under one acquisition should not be assumed to
hold under another. The noise control is what makes this visible ---
without it, the EI column alone reads as a clean positive on both blocks.

Three conclusions follow.

\begin{enumerate}[leftmargin=1.6em,itemsep=0.25em]
  \item \textbf{The null-information gain is an attribution confound.} Under greedy it is
    not significant on either block ($p = 0.51$, $0.57$): we detect no positive
    null-auxiliary advantage once the variance channel is removed. With 32 campaigns
    per cell that is a failure to detect a departure, not a demonstration of
    equivalence; we did not pre-specify an equivalence margin.
  \item \textbf{Naive pooling does not beat the fitted model.} Under greedy,
    \pooled{} on CYP is worse than single-task ($-20.6$ with noise, $-7.5$ with
    the measured panel). An earlier draft concluded the opposite from the EI
    column alone. \pooled{} reported the smallest uncertainty and therefore the
    most exploitative acquisition; it did not use the auxiliary data better.
  \item \textbf{Transfer is real and larger than the EI column suggests.} On JAK
    under greedy, \mt{} gains $+16.2$ points with the measured panel against
    $+3.1$ with noise --- same surrogate, same budget, same history indices,
    only the auxiliary values differ. This is our cleanest evidence that the
    effect is information transfer rather than bookkeeping.
\end{enumerate}


\section{Which number should the model predict?}

One of the largest effects we measure, and the only one requiring no new
machinery. A
campaign returns whatever scalar it was asked to maximise.

\label{sec:selectivity}

Sections above optimise one potency column. Programmes more often fail on
off-target activity than on potency, so we optimise the contrast
$y_{\text{sel}}(m) = \mathrm{pIC}_{50}(m, t_0) - \max_{j \neq 0}
\mathrm{pIC}_{50}(m, t_j)$ on the same blocks.

Every arm observes the full panel on any compound it selects, which is what
selectivity profiling is. Budget is counted in compounds, so no arm sees more
measurements than another.

\paragraph{Selectivity is a different search.} The selectivity and potency
top-10 sets are disjoint on JAK1, JAK2 and TYK2 ($0\%$ overlap) and share $40\%$
on JAK3. Optimising potency does not approximate optimising selectivity.

\begin{table}[ht]
\centering\small
\begin{tabular}{lrr}
\toprule
Arm & JAK ($B = 200$) & CYP ($B = 74$) \\
\midrule
\rnd{}                                   & $21.9\%$ & $21.6\%$ \\
potency (optimise on-target, score selectivity) & $55.9\%$ & $41.6\%$ \\
\textbf{contrast} (single-task GP on $y_{\text{sel}}$) & $\mathbf{95.6\%}$ & $36.9\%$ \\
compose (one sharing model per target, then combine)  & $93.4\%$ & $38.4\%$ \\
\bottomrule
\end{tabular}
\caption{Top-10 \emph{selectivity} recall at $22\%$ of library, 32 campaigns per
cell. Modelling selectivity beats optimising potency by $+39.7$ points on JAK
($p < 10^{-4}$); on CYP no arm separates.}
\label{tab:selectivity}
\end{table}

Two results, pointing opposite ways.

\textbf{Objective specification produced the largest effect in this
experiment}: $+39.7$
points on JAK for changing which scalar is regressed, at identical cost on
identical data.

\textbf{The multi-task surrogate does not earn its place even here.} Selectivity
is where a joint model should have a structural advantage: it models
every task and can express a between-task contrast, which a single-task model
cannot. It ties instead ($-2.2$ and $+1.6$ points, both n.s.). \texttt{compose}
carries four times the rows, and \S\ref{sec:null} shows those extra rows
\emph{flatter} an arm under EI, so the null holds against a favourable bias.

On CYP no arm separates from potency, the same family split that runs through
the report.

\paragraph{Is the target set real, or the assay's reporting floor?} A difference
objective amplifies censoring: if an anti-target value is a reporting floor,
$y(\text{on}) - \max(\text{off})$ is large for a reason unrelated to
selectivity. Our censoring filter removes values held by $>2\%$ of a column, but
a floor just under that threshold survives. We therefore inspect the winners
directly. Across both blocks the binding anti-target values of the top-10 span
$3.70$--$7.82$ --- the real dynamic range --- and are shared by a median of
1--2 compounds, so the optimisation target is genuine chemistry. One
contaminated compound was found and is reported rather than removed (panel
$[4.7, 4.7, 9.7, 4.7]$: three identical values are the signature of ``inactive
at the screening dose''), sitting in JAK3's top-10. CYP3A4 shows $20\%$ of its
top-10 on repeated values and CYP is the more censored block, which is a second
reason to read its null cautiously.

\section{Is a model worth fitting at all?}

Whether a surrogate is worth fitting at all. The comparison that settles it is
against similarity search, not against random.


\begin{table}[ht]
\centering\small
\begin{tabular}{lrrrr}
\toprule
Block & \rnd{} & \simarm{} & \stl{} & \mt{} \\
\midrule
JAK & 6.0\% & 10.5--13.5\% & 27.2\% & \textbf{27.0--32.8\%} \\
CYP & 17.3\% & 25.3--28.7\% & 38.0\% & \textbf{36.2--40.7\%} \\
\bottomrule
\end{tabular}
\caption{Top-10 recall at a 60-assay budget, top-10 withheld from history.
Ranges span five history compositions.}
\label{tab:baselines}
\end{table}

Similarity search is \emph{worse than plain single-task BO} on both blocks, by
$13.7$ points on JAK at its best history composition. Since it uses exactly the
information the multi-task arm uses
and costs nothing to run, this is the comparison that establishes the modelling
earns its place.

\section{Should related targets share a model?}
\label{sec:joint}

Two questions with opposite answers, routinely stated as one.

\subsection{Splitting a fresh budget across targets does not work}

\begin{table}[ht]
\centering\small
\begin{tabular}{llrr}
\toprule
Dataset & Design & Best $\Delta\rho$ vs \stl{} & Targets improving \\
\midrule
DOCKSTRING (21 kinases) & spread  & $-0.056$ & 0 / 21 \\
DOCKSTRING (21 kinases) & blocked & $-0.189$ & 0 / 21 \\
JAK block (measured)    & spread  & $-0.127 \pm 0.014$ & --- \\
CYP block (measured)    & spread  & $-0.045 \pm 0.008$ & --- \\
\bottomrule
\end{tabular}
\caption{budget-splitting at equal cells. Negative everywhere, on simulated and measured
data alike.}
\label{tab:h1split}
\end{table}

At equal cell budget, spreading assays across coupled targets is consistently
\emph{worse} than concentrating them. This is not a weak-model result: the
model's own fit score selects near-maximal task correlation, and a non-deployable
ceiling arm given the auxiliary columns for free gains $+0.31$ to $+0.55$. The
structure is real; paying for it out of the same budget is what fails.

The mechanism is visible in the gap between designs. Measuring a second target
on the \emph{same} compound is nearly a duplicate reading when tasks are
strongly coupled; measuring it on a \emph{different} compound buys coverage of
new chemistry and recovers most of the loss ($-0.189 \to -0.056$). But even the
better design loses: an auxiliary cell is worth slightly less than a primary
one.

\subsection{Borrowing old data on relatives does work, on one kind of panel}

\begin{table}[ht]
\centering\small
\begin{tabular}{lrrr}
\toprule
Dataset & $n = 10$ & $n = 20$ & $n = 40$ \\
\midrule
DOCKSTRING (simulated)                 & $+0.310$ & $+0.233$ & --- \\
ChEMBL, separate libraries (measured)  & $+0.001$ & $-0.020$ & --- \\
JAK block (measured, co-measured)      & $+0.101$ & $+0.086$ & $+0.049$ \\
CYP block (measured, co-measured)      & $+0.256$ & $+0.181$ & $+0.117$ \\
\bottomrule
\end{tabular}
\caption{borrowed-history: $\Delta$ rank correlation over single-task, with $H = 50$ free
historical observations per relative --- one $H$ for every row, so the rows are
comparable. Nothing in the ChEMBL row clears three standard errors
($\pm 0.022$, $\pm 0.011$). The DOCKSTRING and ChEMBL rows are paired runs
(\texttt{-{}-c-paired -{}-c-holdout 0.4}, \texttt{-{}-paired -{}-holdout 0.4});
an earlier version of this table mixed $H = 200$ for those two rows with
$H = 50$ for the blocks, and read a budget of $n = 80$ into the $n = 40$
column. Dashes are budgets the corresponding grid did not run.}
\label{tab:h1cold}
\end{table}

At matched $n$ and $H$ the DOCKSTRING effect is $+0.310$ against $+0.256$
(CYP) and $+0.101$ (JAK) --- between $1.2\times$ and $3\times$ the measured one,
which is why the simulated result cannot be quoted on its own: docking scores
across pockets are nearly one function (fitted $r = 1.00$ at every budget)
because generic ligand size and lipophilicity drive much of the variance.
Measured panels sit at $r \approx 0.6$ and still show a gain.

The ChEMBL null is important and unexplained. Those four targets are separate
kinase \emph{families} with disjoint libraries; the blocks are single families
with shared compounds. We tested chemical overlap as the explanation and
falsified it (\S\ref{sec:falsified}). \textbf{Why transfer fails across
separately assembled libraries remains open}, and is the most important loose
end in this work.

\subsection{What the old data is worth, in assays}

A single-task campaign at budget 90 reaches $37.7\%$ recall; the multi-task
campaign at budget 60 with $H = 50$ per relative reaches $37.3\%$. So $150$ free
auxiliary assays buy what $30$ primary assays buy: in this one campaign-matching calculation, roughly five historical auxiliary
measurements were needed to replace one primary measurement. This is consistent with budget-splitting's
rejection --- the same fact in campaign units --- and it means the result
depends on the history being already paid for. A paper quoting the recall ratio
without this invites a reviewer to derive it.

\section{How much of the library to test}

Below a floor the method costs recall; above a ceiling single-task BO has
already found everything.

\subsection{The budget range where it helps}
\label{sec:window}

\begin{table}[ht]
\centering\small
\begin{tabular}{rrrrcrrc}
\toprule
& \multicolumn{3}{c}{JAK block ($N = 904$)} & &
  \multicolumn{3}{c}{CYP block ($N = 335$)} \\
\cmidrule(lr){2-4}\cmidrule(lr){6-8}
\% library & \stl{} & \mt{} & $\Delta$ (targets up) & &
             \stl{} & \mt{} & $\Delta$ (targets up) \\
\midrule
11\% & 40.0\% & 46.6\% & $+6.6$ \ (2/4) & & 17.5\% & 18.1\% & $+0.6$ \ (2/4) \\
\textbf{22\%} & 68.8\% & 75.3\% & $\mathbf{+6.6}$ \ \textbf{(4/4)} & &
  37.8\% & 52.2\% & $\mathbf{+14.4}$ \ \textbf{(4/4)} \\
\textbf{33\%} & 76.2\% & 83.4\% & $\mathbf{+7.2}$ \ \textbf{(4/4)} & &
  58.1\% & 70.0\% & $\mathbf{+11.9}$ \ \textbf{(4/4)} \\
50\% & 92.8\% & 94.1\% & $+1.2$ \ (2/4) & & 78.1\% & 87.5\% & $+9.4$ \ (4/4) \\
66\% & 97.8\% & 98.1\% & $+0.3$ \ (1/4) & & 89.1\% & 94.1\% & $+5.0$ \ (3/4) \\
\bottomrule
\end{tabular}
\caption{Top-10 recall against fraction of the library screened. Batch 10,
8 repeats, primary top-10 withheld from history. ``Targets up'' counts how many
of the four held-out targets individually improved.}
\label{tab:window}
\end{table}

Table~\ref{tab:window} is the central practical result and it replicates across
two protein families on the library-fraction axis:

\begin{enumerate}[leftmargin=1.6em,itemsep=0.2em]
  \item \textbf{Below ${\sim}5\%$ the method costs you.} At a 20-assay budget
    the gain is $-0.7$ (JAK) and $-4.3$ (CYP) recall points.
  \item \textbf{At $22$--$33\%$ it is both effective and reliable}, $+6.6$ to
    $+14.4$ points, with \emph{all four} held-out targets improving on both
    blocks.
  \item \textbf{Beyond $50\%$ it decays} as the single-task baseline saturates
    (JAK \stl{} reaches $92.8\%$); there is nothing left to win.
\end{enumerate}

Point 2 matters more than the effect size. At small budgets the per-target
outcome is a lottery: at a 60-assay JAK budget, one target gained 15 recall
points and another lost 15, pooling to a meaningless ``$+3.7\%$''. A
practitioner screening one cold target gets one draw, not an average. The
unanimity at $22$--$33\%$ is what makes a leave-one-target-out self-check
(\S\ref{sec:guidance}) able to license a decision about an unmeasured target.

\subsection{A share of the library, or a number of assays?}
\label{sec:scaling}

Whether the window is a fraction or an assay count decides whether the method is
usable: screening a third of a catalogue of millions is impossible, spending
$300$ assays is routine. We varied library size and absolute budget
independently on the JAK block ($N \in \{904, 600, 400\}$ by random subsample,
budget $\in \{100, 150, 200, 300\}$), giving matched-fraction pairs that differ
in budget and matched-budget columns that differ in fraction.

A first pass held $H = 50$ per relative and appeared to show absolute budget
governing ($R^2 = 0.30$ against $0.18$). That was an artefact of the design:
with $H$ fixed, shrinking the library raises the history's coverage of it from
$17\%$ to $38\%$, favouring the cells that also have the highest budget
fraction. Holding coverage at $16.6\%$ reverses the ranking.

\begin{table}[ht]
\centering\small
\begin{tabular}{lrr}
\toprule
Predictor of \mt{} $-$ \stl{} & $R^2$, $H$ absolute & $R^2$, coverage held fixed \\
\midrule
Fraction of library screened & $0.18$ & $\mathbf{0.47}$ \\
Absolute assay count         & $0.30$ & $0.16$ \\
Single-task headroom $(1 - \text{recall})$ & $0.31$ & $0.43$ \\
\bottomrule
\end{tabular}
\caption{12 cells, 32 campaigns each. Controlling history coverage swaps the
two candidate invariants.}
\label{tab:scaling}
\end{table}

With coverage controlled, matched-fraction pairs agree (differences $-5.0$,
$-0.9$, $+0.9$, $+3.1$ points, none significant) and fraction is the better
predictor. The window is a fraction, not an assay count, and the history must be
a fraction too: both scale with library size, so a deployment must satisfy both.
Fraction and headroom are near-collinear ($0.47$ against $0.43$) and we cannot
separate them.

$R^2 = 0.47$ on 12 cells with $2.5$-point standard errors leaves most of the
variance unexplained, and the uncontrolled analysis pointed the other way.
Resolving the window's edges would need about four times the repeats.

\subsection{Starting away from the answer}
\label{sec:scaffold}

A random initial design can seed a campaign with a direct analogue of the
answer. On these panels that is a live concern: the true top-10 spans only
$6$--$9$ distinct core scaffolds out of $357$ in the JAK block, so the
compounds being searched for are concentrated in a few congeneric series.

The virtual-screening literature handles this with a scaffold split of the
train/test partition. A campaign has no fixed training set --- it chooses one
--- so we apply the same rule to the \emph{initial design}: the seed compounds
are drawn only from scaffolds absent from the true top-$k$, and the campaign may
then select anything. It must therefore cross scaffolds to succeed, which is the
retrospective analogue of screening a catalogue in which one holds no series.

\begin{table}[ht]
\centering\small
\begin{tabular}{lrrrrrr}
\toprule
 & & \multicolumn{2}{c}{random seed} & \multicolumn{2}{c}{scaffold-split seed} & \\
\cmidrule(lr){3-4}\cmidrule(lr){5-6}
Block & \% lib & \stl{} & $\Delta$ (up) & \stl{} & $\Delta$ (up) & $p$ \\
\midrule
JAK & $22\%$ & $68.8\%$ & $+6.6$ (4/4) & $68.4\%$ & $+5.9$ (3/4) & $0.84$ \\
JAK & $33\%$ & $76.2\%$ & $+7.2$ (4/4) & $79.4\%$ & $+3.1$ (3/4) & $0.14$ \\
CYP & $22\%$ & $37.8\%$ & $+14.4$ (4/4) & $41.2\%$ & $+12.5$ (3/4) & $0.69$ \\
CYP & $33\%$ & $58.1\%$ & $+11.9$ (4/4) & $62.8\%$ & $+9.4$ (3/4) & $0.57$ \\
\bottomrule
\end{tabular}
\caption{Top-10 recall with the initial design drawn at random, and restricted
to scaffolds absent from the true top-10. 32 campaigns per cell; $p$ compares
the two $\Delta$ distributions.}
\label{tab:scaffold}
\end{table}

The effect size survives. \mt{} still beats \stl{} in all four cells, the gain
shrinks in all four, and no shrinkage is significant ($p = 0.14$ to $0.84$).
The largest drop, $+7.2 \to +3.1$ on JAK at $33\%$, is also the cell where the
single-task baseline rises most ($76.2 \to 79.4\%$), so most of it is the
baseline improving rather than the multi-task arm degrading.

\textbf{Unanimity does not survive.} Every cell falls from $4/4$ held-out
targets improving to $3/4$. That is the property \S\ref{sec:window} identifies
as the reason the operating window is usable at all: a leave-one-target-out
self-check can only license a decision about an unmeasured target if its answer
does not depend on which target was held out. Under a scaffold-split seed it
does, in all four cells we tested.

We therefore weaken the claim. The multi-task gain is not an artefact of seeding
the campaign with analogues of the answer --- that was the concern a scaffold
split tests, and it is not supported. But the unanimity that made the window a
recommendation rather than an average is contingent on a random seed, and a
practitioner holding no series against the target is in the $3/4$ regime, not
the $4/4$ one.

One property of these panels explains why the split bites at all: the true
top-10 spans only $6$--$9$ core scaffolds out of $357$ in the JAK
block. Withholding those scaffolds removes $4$--$19\%$ of the library from the
seed pool. On a library with less series structure the split would be closer to
a no-op, and on a catalogue it would be stricter.

\section{Where to put the other requirements}

Where a cheap descriptor enters the loop matters more than which descriptor it
is.

\label{sec:filter}

Sections above optimise raw potency, so nothing prevents the method selecting a
compound that binds but is undevelopable. Adding a drug-likeness \emph{term} is
not the repair --- that is what ruins the published DOCKSTRING objectives
(\S\ref{sec:data-quality}). Applied as a library \emph{admission filter}
($\mathrm{QED} \ge 0.5$) the score stays entirely assayed and the zero-assay
floor stays at chance.

The filter is unbiased with respect to potency: retention of each target's
top-10 (4--8 of 10) tracks the overall $63\%$ rate, and the top-10's mean QED
($0.51$--$0.64$) matches the library's ($0.54$). JAK at $33\%$ of the library:

\begin{table}[ht]
\centering\small
\begin{tabular}{lrrc}
\toprule
Variant & $N$ & $\Delta$recall & Targets improving \\
\midrule
Full clean block           & 904 & $+7.2$ & \textbf{4/4} \\
$\mathrm{QED} \ge 0.5$     & 566 & $+6.9$ & 3/4 \\
Random cut (size control)  & 566 & $+2.5$ & 1/4 \\
\bottomrule
\end{tabular}
\caption{The drug-like subset preserves the effect; a random cut of equal size
does not. The loss of unanimity is a library-size effect, not the filter.}
\label{tab:qed}
\end{table}

We recommend keeping the filter: it costs ${\sim}0.3$ recall points, removes the
objection that the method ignores developability, and never enters the score.
Note that \textbf{the unanimity result of Table~\ref{tab:window} is demonstrated
at $N \approx 900$}; it should not be assumed to hold at arbitrary library size.

\paragraph{Why not drug-likeness as a second objective?} We implemented that
too --- a standard multi-objective rule over (potency, drug-likeness), reweighted each round
per round, Pareto-front recall as the metric --- expecting a Pareto treatment to
avoid the single-score failure of \S\ref{sec:data-quality}. It does not. A
\emph{free} continuous objective puts part of the front on sale: the maximum-QED
compound is non-dominated by construction, so ranking by QED alone and measuring
nothing recovers $20$--$50\%$ of the front. The fronts are also tiny (2--10
compounds of 335--904, so one compound is $10$--$50\%$ of the metric) and
unrepresentative (median QED on the front $0.72$--$0.91$ against a library
median of $0.54$--$0.58$). At budget 100 on JAK, the zero-assay \textsc{qed-only}
arm reaches $50.0\%$ front recall against \mt{}'s $50.0\%$ and \stl{}'s $51.7\%$.

The distinction that matters is \emph{what each construction changes}: a score
term changes the score, a Pareto objective changes the success set, and a filter
changes only which compounds are eligible. Only the last leaves the objective
fully assayed and the zero-assay floor at chance.

\paragraph{Where a free descriptor belongs.} A fourth option is to optimise
potency over the whole library and apply drug-likeness only when scoring, since
QED is free for any SMILES. It is not obviously worse: a potent but
undevelopable compound still teaches the surrogate about the SAR, which a
pre-filter discards. We compared all of them at \emph{identical success set}
(top-10 potency among the 566 drug-like) and \emph{identical absolute budget};
only where the budget may be spent, and what the free history covers, differ.

\begin{table}[ht]
\centering\small
\begin{tabular}{rlrrr}
\toprule
Budget & Design & \stl{} & \mt{} & $\Delta$recall \\
\midrule
62  & post-hoc: search all 904            & 19.4\% & 14.7\% & $-4.7$ \\
    & pre-filter: library cut to 566      & 27.5\% & 31.9\% & $+4.4$ \\
    & \textbf{restricted: spend 566, history 904} & 27.5\% & \textbf{37.5\%} & $\mathbf{+10.0}$ \\
\midrule
125 & post-hoc                            & 42.2\% & 45.3\% & $+3.1$ \\
    & pre-filter                          & 52.5\% & 58.4\% & $+5.9$ \\
    & \textbf{restricted}                 & 52.5\% & \textbf{59.7\%} & $+7.2$ \\
\midrule
187 & post-hoc                            & 51.9\% & 57.5\% & $+5.6$ \\
    & \textbf{pre-filter}                 & 66.6\% & \textbf{73.4\%} & $+6.9$ \\
    & restricted                          & 66.6\% & 70.0\% & $+3.4$ \\
\bottomrule
\end{tabular}
\caption{Post-hoc scoring costs $13$--$17$ recall points: ${\sim}37\%$ of the
library cannot count as a success, so ${\sim}37\%$ of the budget buys compounds
that will be discarded. Confining the \emph{budget} while letting the free
history span the whole library is best at the budgets where the method matters;
at 187 the two converge within the noise of 8 repeats. That \stl{} is identical
under \textsc{pre-filter} and \textsc{restricted} ($27.5/52.5/66.6$) is the
implementation check --- single-task uses no history, so the two must coincide
for it.}
\label{tab:where}
\end{table}

\begin{quote}
\textbf{General lesson.} A free descriptor has four possible roles and they are
not interchangeable. As a term in the score it dominates the ranking
(\S\ref{sec:data-quality}, 93\%). As a Pareto objective it owns a region of the
front and gives away part of the answer (20--50\% of a 2--10 compound front).
Applied only when scoring, it wastes the budget it should have steered.
\textbf{A free quantity should constrain what you pay for, and nothing else.}
Correspondingly, any multi-objective molecular BO benchmark should report what a
zero-assay ranking on each free objective recovers, and treat that as its floor.
\end{quote}

\subsection{Sensitivity to the censoring threshold}
\label{sec:pilefrac}

Our censoring filter drops any compound holding a value that more than $2\%$ of
a column shares. It does a great deal of work --- $57.8\%$ of the JAK block and
$22.8\%$ of the CYP block survive it --- and $2\%$ is a heuristic. We sweep it,
holding the budget at $22\%$ of whatever library survives so that every row sits
at the same point of the coverage window.

\begin{table}[ht]
\centering\small
\begin{tabular}{lrrrrrrr}
\toprule
Block & threshold & $N$ & budget & \stl{} & \mt{} & $\Delta$ & targets up \\
\midrule
JAK & $0.5\%$ & $417$ & $92$  & $50.3\%$ & $65.6\%$ & $+15.3^{*}$ & 4/4 \\
JAK & $1\%$   & $871$ & $192$ & $65.9\%$ & $72.8\%$ & $+6.9^{*}$  & 2/4 \\
JAK & $2\%$   & $904$ & $199$ & $68.8\%$ & $75.3\%$ & $+6.6^{*}$  & 4/4 \\
JAK & $3\%$   & $919$ & $202$ & $66.6\%$ & $67.8\%$ & $+1.2$      & 3/4 \\
JAK & $5\%$   & $927$ & $204$ & $64.1\%$ & $70.3\%$ & $+6.2^{*}$  & 3/4 \\
JAK & $10\%$  & $982$ & $216$ & $66.2\%$ & $69.7\%$ & $+3.4$      & 2/4 \\
\midrule
CYP & $0.5\%$ & $201$   & $44$  & $36.2\%$ & $49.7\%$ & $+13.4^{*}$ & 2/4 \\
CYP & $1\%$   & $285$   & $63$  & $38.8\%$ & $49.4\%$ & $+10.6^{*}$ & 2/4 \\
CYP & $2\%$   & $335$   & $74$  & $37.8\%$ & $52.2\%$ & $+14.4^{*}$ & 4/4 \\
CYP & $3\%$   & $350$   & $77$  & $39.7\%$ & $47.2\%$ & $+7.5$      & 4/4 \\
CYP & $5\%$   & $526$   & $116$ & $48.8\%$ & $61.9\%$ & $+13.1^{*}$ & 4/4 \\
CYP & $10\%$  & $1{,}036$ & $228$ & $66.9\%$ & $72.5\%$ & $+5.6$    & 3/4 \\
\bottomrule
\end{tabular}
\caption{Censoring-threshold sensitivity, budget held at $22\%$ of the surviving
library. 32 campaigns per cell; $^{*}$ marks $p<0.05$ on a paired test, with the
multiplicity caveat of \S\ref{sec:stats}.}
\label{tab:pilefrac}
\end{table}

\textbf{The sign is robust: \mt{} is ahead in all twelve cells.} The magnitude
is not. On JAK it ranges from $+1.2$ to $+15.3$, and eight of twelve cells reach
$p<0.05$ while four do not.

The most informative comparison is between adjacent thresholds. From $1\%$ to
$5\%$ the JAK library changes by $6\%$ in size ($871$ to $927$) and the budget by
$12$ compounds, yet the measured gain moves $+6.9$, $+6.6$, $+1.2$, $+6.2$.

These are \emph{different datasets}, not repeated runs of one: compounds enter
and leave, the budget shifts, and the success set can change with them. The
roughly five-point spread is therefore a measure of sensitivity to an analysis
choice, not an estimate of campaign-seed variability. What makes it worth
reporting is that it \emph{exceeds} the within-configuration Monte Carlo
standard error of $\pm 2.5$ points: for effects of this size, how the data are
filtered matters at least as much as which seeds are drawn.

We use this as a heuristic robustness scale rather than a statistical error
estimate. Effects below roughly five points in this report --- including some we
mark significant --- should be read as provisional until shown stable across
analysis choices, not merely across seeds.

Two further observations. The aggressive $0.5\%$ threshold produces the largest
gains on both panels ($+15.3$, $+13.4$) on the smallest libraries, consistent
with \S\ref{sec:scaling}: coverage is held fixed but the absolute budget falls
to $92$ and $44$. And at $10\%$ the filter barely filters --- $982$ of $1{,}563$
on JAK and $1{,}036$ of $1{,}472$ on CYP --- so censored values are substantially
back in the library, and the gain is smallest and not significant on either
panel. That is the direction the filter was introduced to correct, and it
supports keeping a filter without pinning the exact threshold.

\section{Things we expected that turned out false}
\label{sec:falsified}

We report these because each was a plausible mechanism that a reader might
otherwise propose, and because the effort spent excluding them is part of what
the positive result rests on.

\begin{description}[leftmargin=2em,style=nextline,itemsep=0.35em]
  \item[A full task-covariance matrix should beat one shared $r$.]
    Motivated by the JAK family being unevenly coupled ($0.24$ to $0.71$) where
    CYP is uniform ($0.52$ to $0.74$). We fitted a free correlation matrix by
    coordinate ascent on the exact likelihood with PSD rejection. It \emph{loses}
    on both blocks ($-0.003$ to $-0.029$), largest where there is most data to
    overfit. The fitted matrix does recover the heterogeneity (spread $0.49$ vs
    empirical $0.47$); six correlations estimated from 10--40 primary
    observations simply cost more in variance than uneven coupling costs in
    bias. \textbf{One shared $r$ is the right amount of structure at these
    budgets.}
  \item[Chemical overlap should explain the ChEMBL null.]
    Real per-target libraries reach only $0.53$ of the primary's own
    nearest-neighbour similarity, against $1.0$ for a co-measured block, which
    made overlap the natural explanation. We made it a controlled variable and
    swept it. The campaign gain is \emph{flat} from reach $0.96$ to $0.44$
    ($+0.084 \to +0.077$ on JAK, SEM $0.023$) --- and $0.44$ is \emph{below} the
    real-library value. Overlap is not the mechanism.
  \item[The regression gain should predict the campaign gain.]
    The rank correlation gain shrinks with budget ($+0.20$ at $n=10$ to $+0.06$ at
    $n=80$), suggesting an early-campaign effect. The campaign gain does the
    opposite: negative at 20 assays, peaking at $22$--$33\%$ of the library. A
    regression curve does not predict the campaign curve, even in sign.
  \item[Coupling heterogeneity should explain when fitting $r$ beats pooling.]
    \pooled{} loses badly on JAK and ties on CYP, and heterogeneity predicted
    this. But it also predicted the free-$B$ result, which failed, and under the
    stricter control (top-10 withheld) the JAK/CYP split does not survive.
    The mechanism is withdrawn; \S\ref{sec:null} supplies a different one.
  \item[Naive pooling should outperform the fitted joint modelling.]
    \emph{This was our own conclusion in an earlier draft}, from EI columns in
    which \pooled{} beat \mt{} on CYP by up to $9.1$ points at $p = 0.0004$. It
    is wrong. \pooled{} reports the smallest uncertainty of the three
    arms, which makes EI most exploitative; under a variance-free acquisition it
    is \emph{worse} than single-task on CYP. We had run the pooling control that
    standing practice asks for, but had not asked whether the control was
    comparable on the axis the acquisition actually reads. It was not.
  \item[Absolute assay count should govern the operating window.]
    Supported at $R^2 = 0.30$ against $0.18$ for library fraction --- until the
    history's coverage of the library was held fixed, at which point the two
    swap ($0.16$ against $0.47$). See \S\ref{sec:scaling}.
  \item[The gain should come from seeding the campaign near the answer.]
    The top-10 of these panels sits in $6$--$9$ scaffolds out of $357$, so a
    random initial design can hand a campaign a direct analogue of what it is
    looking for. Restricting the seed to scaffolds absent from the top-10
    (\S\ref{sec:scaffold}) shrinks the gain in all four cells and in none of
    them significantly ($p = 0.14$ to $0.84$). The effect is not an artefact of
    the seed. What the split does remove is the \emph{unanimity}, which falls
    from $4/4$ to $3/4$ held-out targets in every cell.
  \item[A computed auxiliary should substitute for a measured panel.]
    The natural route to deployment, since it would remove the
    co-measured-panel requirement. Degrading the auxiliary columns to a
    controlled rank correlation $\rho$ against their own values puts the threshold at
    $\rho \approx 0.45$--$0.6$, above the $0.2$--$0.4$ that docking scores reach
    against measured pIC$_{50}$. We therefore did not build a docking cache.
    The threshold is provisional: it was measured under EI, against the
    contaminated baseline of \S\ref{sec:null}, which inflates the low-$\rho$ end
    and so makes the threshold look lower than it is.
\end{description}

\section{What each choice is worth}
\label{sec:worth}

Effect sizes of five campaign decisions, in top-10 recall points at a $22\%$
library budget. The rows are not a ranking: two are panel-dependent and one is
negative on a panel, as the text explains.

\begin{table}[ht]
\centering\small
\begin{tabular}{llrl}
\toprule
Decision & Choice that pays & Worth & Conditions \\
\midrule
What to regress & the contrast you actually want & $\mathbf{+40}$ &
  none --- same data, same cost \\
Fit a surrogate & GP + EI over similarity search & $+42$ / $-3$ &
  panel-dependent; verify \\
Model jointly & only with free co-measured history & $+7$ to $+14$ &
  three conditions, all required \\
Budget & $22$--$33\%$ of the library & --- &
  below ${\sim}5\%$ it is negative \\
Secondary criteria & admission filter, not score term & $+13$ to $+17$ &
  vs.\ post-hoc scoring \\
\bottomrule
\end{tabular}
\caption{The choice of what to optimise produced one of the largest effects we measured.}
\label{tab:worth}
\end{table}

Two entries are panel-dependent, and the ordering should not be read as a
ranking. Fitting a model at all is worth $+42$ points on JAK and $-3$ on CYP,
where similarity search ties single-target BO --- a larger maximum than the
objective-choice row, on one panel, and a negative one on the other. Only the
data-sharing arm beats the model-free heuristic on both panels. Joint modelling
requires all three of its conditions at once --- history already paid for,
co-measured on the same compounds, and a library small enough that $22$--$33\%$
is affordable --- and fails when the quantity of interest is already computable
from data in hand.

Figures in the joint-modelling row should be read against the
\emph{real-versus-null auxiliary contrast} of \S\ref{sec:null}: the difference
between an arm given measured auxiliary columns and the same arm given noise
columns, with everything else held fixed. We do not subtract the two and call
the remainder ``the transfer effect'' --- that would assume the informational
and uncertainty-scale effects add, and they may interact. The contrast isolates
the consequence of changing the auxiliary labels from uninformative to measured,
which requires no such assumption.

\section{Limits, and what we would run next}
\label{sec:limitations}

\begin{itemize}[leftmargin=1.6em,itemsep=0.2em]
  \item \textbf{Retrospective.} Every ``assay'' is a lookup from a published
    table. No compound was synthesised or prospectively tested.
  \item \textbf{Single objective.} We optimise one potency column. Real
    programmes optimise potency, selectivity, solubility and metabolic stability
    jointly, and fail on the latter three. Our one attempt at a published
    composite objective found it $93\%$ determined by a free descriptor.
  \item \textbf{Two blocks, eight targets.} Both are co-measured panels from one
    source.
  \item \textbf{Fixed library.} The method selects from compounds that already
    exist; it does not design or generate.
  \item \textbf{Transfer failed in our separately assembled cross-family
    ChEMBL setting and the cause remains unresolved.} We have excluded chemical
    overlap, but that setting differs from the successful panels in several
    respects at once (\S\ref{sec:guidance}), so we cannot attribute the failure
    to any one of them.
  \item \textbf{Batch size is a confound across sweeps.} Batch 5 and batch 10
    runs are not directly comparable; comparisons are valid within a sweep only.
  \item \textbf{Arm comparisons are acquisition-dependent.} \S\ref{sec:null}
    shows that the measured transfer effect changes sign-of-significance across
    EI, greedy and Thompson sampling, and that a null auxiliary gains recall
    under EI alone. Every effect size here should be read against its matched
    $\rho = 0$ control \emph{under the same acquisition}. A variance-calibrated
    acquisition would be the correct fix and is not implemented.
  \item \textbf{One surrogate family.} We use an exact GP with a Tanimoto
    kernel. Recent cross-domain work reports neural surrogates with Thompson
    sampling outperforming GPs across 61 tasks, and our Thompson-sampling
    results suggest the interaction between surrogate capacity and acquisition
    is where the remaining uncertainty lies.
  \item \textbf{Unanimity is demonstrated at $N \approx 900$ and $N \approx
    335$.} A random cut of the JAK block to $566$ compounds reduces the gain from
    $+7.2$ to $+2.5$ and unanimity from $4/4$ to $1/4$. The window is a claim
    about these blocks at these sizes.
  \item \textbf{Our effects are sensitive to analysis choices, not only to
    seeds.} Four nearby censoring thresholds on JAK give gains of $+6.9$,
    $+6.6$, $+1.2$ and $+6.2$ (\S\ref{sec:pilefrac}). Those are different
    datasets rather than repeated runs, so the roughly five-point spread
    measures sensitivity to an analysis choice --- and it exceeds the
    $\pm 2.5$ point Monte Carlo standard error within a configuration. We use
    five points as a heuristic robustness scale, not a statistical error
    estimate: effects below it should be read as provisional until shown
    stable across analysis choices, not merely across seeds.
  \item \textbf{Two conclusions reversed under a control run after the fact},
    and in both cases the uncontrolled reading was the confident one. We regard
    the remaining conclusions as provisional to the same degree and state which
    control each is waiting on.
  \item \textbf{Unanimity is contingent on the initial design.} It holds $4/4$
    with a random seed and $3/4$ with a scaffold-split seed, in every cell we
    tested (\S\ref{sec:scaffold}). The window remains a positive-expectation
    recommendation; it is not a per-target guarantee.
  \item \textbf{Scaffold splitting is applied to the initial design only.} We do
    not scaffold-split a train/test partition, because a campaign has no fixed
    training set. A benchmark that reports a regression $\rho$ on these panels
    should do so.
\end{itemize}

\section{What to do with this}
\label{sec:guidance}

For a library of known compounds not yet assayed. The first item produced the largest effect in our selectivity experiment and
needs no data sharing at all.

\begin{enumerate}[leftmargin=1.6em,itemsep=0.2em]
  \item \textbf{Regress the objective you actually want.} If you want selective
    compounds, regress the selectivity contrast, not potency: $+40$ recall
    points on JAK at identical cost (\S\ref{sec:selectivity}), and the potency
    and selectivity top-10 sets are disjoint on three of four targets. A plain
    single-task GP on the contrast is sufficient --- the joint model
    adds nothing here.
  \item \textbf{Audit a contrast objective for censoring before trusting it.}
    A difference amplifies reporting floors; inspect the binding anti-target
    values of your top-$k$ and confirm they are not piled on a limit of
    detection.
\end{enumerate}

For the multi-task transfer result specifically:

\begin{enumerate}[leftmargin=1.6em,itemsep=0.2em]
  \item \textbf{On these panels the gains were largest and most consistent
    around $22$--$33\%$ library coverage.} We have two blocks and eight
    targets, and \S\ref{sec:scaling} resolves the window's edges only weakly
    ($R^2 = 0.47$ on 12 cells), so this is where our measurements peaked rather
    than an estimated constant. Below ${\sim}11\%$ the
    per-target outcome is a lottery and lower still it hurts; beyond $50\%$
    single-task BO already finds almost everything. The invariant is the
    \emph{fraction}, not the assay count (\S\ref{sec:scaling}), which is what
    restricts this to focused libraries of hundreds to low thousands rather than
    catalogue-scale screening.
  \item \textbf{Expect to need history on a comparable fraction.} Our runs hold
    free history on ${\sim}17\%$ of the library, and the benefit does not
    survive holding $H$ fixed while the library grows. Both quantities scale.
  \item \textbf{Use only history you already own.} Auxiliary assays are worth
    roughly one fifth of a primary measurement in our campaign-matching
    calculation. These results do not justify acquiring auxiliary assays solely
    for this purpose. A computed
    substitute would need rank correlation ${\gtrsim}0.45$ against the auxiliary it
    replaces, above what docking reaches.
  \item \textbf{Fit the task correlation; do not pool.} Naive pooling looks
    competitive under EI for a reason unrelated to the auxiliary data
    (\S\ref{sec:null}) and is worse than single-task once that is removed.
  \item \textbf{Evidence of benefit here is limited to the co-measured
    within-family panels.} Our separately assembled per-target libraries showed
    none, but the two settings differ in co-measurement, family relatedness,
    source, assay context and library assembly at once, so we have not isolated
    co-measurement as a necessary condition.
  \item \textbf{Do not substitute similarity search} --- it is worse than plain
    single-task BO.
  \item \textbf{Expect $3/4$, not $4/4$, if you hold no series against the
    target.} Unanimity across held-out targets is measured with a random
    initial design; under a scaffold-split seed one of four targets fails in
    every cell we tested.
  \item \textbf{Run the leave-one-target-out self-check first.} Hold out each
    target you already have, use the rest as history, and confirm multi-task
    beats single-task on all of them at your intended budget. The check is only
    informative in the $22$--$33\%$ window, where we observe unanimity; at
    smaller budgets it gives a different answer depending on which target is
    held out.
\end{enumerate}

\subsection{A baseline set for anyone repeating this}

We propose the following as the minimum comparison set for any multi-task
molecular BO result, because each of them changed a conclusion of ours at least
once:

\begin{enumerate}[leftmargin=1.6em,itemsep=0.2em]
  \item \textbf{Random selection} --- the floor.
  \item \textbf{Similarity search} to known actives --- the free, model-free
    alternative a chemist would actually use.
  \item \textbf{Single-task GP BO at equal \emph{cell} budget} --- not equal
    compounds.
  \item \textbf{Single-task GP BO at a \emph{higher} budget}, to report the
    exchange rate between auxiliary and primary assays.
  \item \textbf{Pooled multi-task} ($r = 1$), to separate task modelling from
    extra rows --- and compared on an axis the acquisition does not distort;
    see the next item.
  \item \textbf{A noise auxiliary} ($\rho = 0$): the multi-task arm
    given auxiliary columns that are pure noise. It should show no systematic advantage in expectation. This is the only check we found that separates information transfer
    from acquisition calibration, it costs one extra run, and it invalidated a
    conclusion of ours that five other controls had passed. We regard it as the
    single most important item on this list.
  \item \textbf{A free-descriptor floor} for any composite objective, and a
    censoring screen for any dense block. For a multi-objective benchmark, report
    what a zero-assay ranking on each \emph{free} objective recovers of the
    Pareto front, and the front's size --- a front of ten compounds cannot
    resolve a method.
\end{enumerate}

Results should be reported per held-out target, not only pooled, with the spread
stated; and with the budget expressed as a fraction of library size. Effect
sizes should be reported alongside their matched null-auxiliary control under
the same acquisition.

\subsection{What we would run next}
\label{sec:next}

Objective variants come first because they require no additional surrogate
machinery and produced a large effect in the JAK selectivity experiment; the
prospective experiment comes last because it changes the category of the work.

\paragraph{Objectives a programme actually specifies.}

\S\ref{sec:selectivity} tests one selectivity objective, a margin against the
worst off-target. Real projects specify other shapes, each a small change to the
same loop:

\begin{itemize}[leftmargin=1.6em,itemsep=0.2em]
  \item a selectivity \emph{window}: a threshold rather than a margin, so
    additional selectivity past the requirement earns nothing;
  \item a \emph{weighted} anti-target panel, because off-targets differ in how
    much they matter;
  \item potency \emph{under a selectivity constraint}, which is the actual
    go/no-go shape --- and which combines with the admission-filter result of
    \S\ref{sec:filter} rather than competing with it.
\end{itemize}

On the evidence of Table~\ref{tab:worth} these should matter more than further
work on the surrogate.

\paragraph{Escaping the retrospective frame.}

Every assay here is a lookup from a table of existing values, on libraries
assembled because the compounds were worth testing. The method never sees the
table --- the runtime audit enforces that --- but the libraries are enriched in
actives, so absolute recall will not transfer to a purchasable catalogue.

\paragraph{A blinded time-split} is the cheapest attack: history before a
deposit-date cutoff, evaluation after it. We note for anyone attempting it that
it is \emph{not} runnable from the data as we cached it --- our ChEMBL extracts
carry no deposit date and our fetcher never requested one. The re-pull is a few
hours and the experiment is free thereafter.

\paragraph{A compute oracle is the decisive step.} The limitations above share
a cause: precomputed values, so the library is finite, enriched and already
known. A structure-based affinity predictor removes all three. The candidate set
becomes any purchasable compound, values are computed on demand, and no value
exists until the campaign pays for it. It is also where the window of
\S\ref{sec:scaling} --- a \emph{fraction} of the library --- meets a library
whose size the experimenter chooses.

The oracle interface is agnostic to whether a value is looked up or computed,
so a co-folding affinity predictor can be substituted directly. We do not report
such an evaluation here: published affinity predictors are trained on the same
public binding data these panels are drawn from, and the resulting train--test
leakage makes their accuracy on public benchmarks uninterpretable. Evaluating
one would require a temporal or scaffold holdout constructed against that
model's training set.

\paragraph{Wet-lab validation} remains the only thing that changes the
category of the result: one cold target with a family panel, a few hundred
assays split between arms, and a threshold-hit metric.

\paragraph{Methodological gaps we are leaving open.}

\begin{description}[leftmargin=2em,style=nextline,itemsep=0.3em]
  \item[A variance-calibrated acquisition.] Arms whose model uncertainty
    differ in scale are not comparable under EI (\S\ref{sec:null}), and our
    diagnostic --- ranking on the mean alone --- is a worse optimiser than the
    setting practitioners use. Calibrating variance across arms would let the
    effect sizes be stated at the operating point rather than at a diagnostic
    one, and would settle the auxiliary-quality threshold.
  \item[Why cross-family transfer fails.] Chemical overlap is excluded. Target
    relatedness is the surviving candidate and is untested.
  \item[A predictor of the per-target outcome.] Unanimity at $22$--$33\%$ makes
    the method usable, but nothing predicts \emph{which} targets gain most.
\end{description}

\section*{Reproducibility}
\label{sec:repro}

All results come from deterministic, seeded runs --- the same command twice
gives byte-identical output.

All five campaign and regression instruments take \texttt{--audit}, which checks
at runtime the invariants the comparisons rest on: that arms start from the same
initial design and see the same history, that each spends exactly its budget
with no repeated pick, that no selection escapes an eligibility restriction,
that no top-$k$ compound leaks into the history when it is meant to be
withheld, and that no auxiliary column is the primary target's own values.
Every configuration behind a number in this report passes, and each run records
its own audit verdict inside its output file, so a table can be checked for
provenance without the terminal that produced it. The checks were validated by
deliberate injection --- perturbing one arm's RNG, adding training rows to the
evaluation set, and inflating a history draw by one row all produced the
expected failures --- because a check that has never fired is not evidence.

The audit also found a defect it was written to find. The two instruments behind
the cold-target results redrew the primary's train/test split for every value of
$H$, so the $H = 0$ control and the arms differenced against it were fitted on
different rows and scored on different rows, in every cell. The gain remains
unbiased, but the comparison was never paired. Re-run paired, the positive
result holds with effect sizes reduced in 15 of 18 cells, and the negative
result on separately assembled libraries holds with its standard errors reduced
by ${\sim}40\%$ --- a null measured on a tighter instrument than the one that
first reported it.

Key scripts:
\texttt{dockstring\_audit.py} (dataset provenance and the QED finding),
\texttt{transfer\_headroom.py} (budget-splitting, ceilings),
\texttt{cold\_target\_chembl.py} (borrowed-history on separate libraries),
\texttt{dense\_block\_h1.py} (dense blocks, censoring filter, free $B$),
\texttt{dense\_block\_campaign.py} (campaigns, baselines, budget and reach
sweeps). Detailed logs, including the intermediate results and corrections that
led to the conclusions above --- including the two conclusions that reversed
under a later control --- are in \texttt{results/FINDINGS.md} (\S\S14--32).

ibliographystyle{plain}
ibliography{refs}

% The citations below were never resolvable: this document cited seven
% keys and declared no bibliography, so every one would have compiled to
% a '[?]'. refs.bib carries all of them.
\bibliographystyle{plain}
\bibliography{refs}

\end{document}
